%! Author = chtompki
%! Date = 1/14/26
% Example LaTeX document for GP111 - note % sign indicates a comment
\documentclass[12pt]{amsart}
% Default margins are too wide all the way around. I reset them here
\setlength{\hoffset}{-.75in}
\setlength{\textwidth}{6.5in}
\setlength{\voffset}{-.5in}
\setlength{\textheight}{9.0in}
\usepackage{amsmath}
\usepackage{amssymb}
\usepackage{amsthm}
\usepackage{pgfplots}
\usepackage{enumerate}
\usepackage{tikz}
\usetikzlibrary{shadows,arrows,arrows.meta,shapes,positioning,calc}% Required for the Circle and Triangle arrow tips
\usepackage{setspace}
\usepackage{siunitx}
\usepackage{enumitem}

\theoremstyle{conjecture}
\newtheorem{conjecture}{Conjecture}

\theoremstyle{definition}
\newtheorem{definition}{Definition}

\theoremstyle{question}
\newtheorem{question}{Question}

\theoremstyle{theorem}
\newtheorem{theorem}{Theorem}

\newenvironment{solution}
{\begin{proof}[Solution]}
{\end{proof}}

\title{Graph Theory:\\Homework 1\\ Book Section 1.1.}
\author{Rob Tompkins}
\date{\today}


\begin{document}

    \maketitle

    % Problem 1
    \textbf{1.} Consider the following four families of graphs:
    \[
        A=\{\text{paths}\}, \qquad
        B=\{\text{cycles}\}, \qquad
        C=\{\text{complete graphs}\}, \qquad
        D=\{\text{bipartite graphs}\}.
    \]
    For each pair of these families, determine all isomorphism classes of
    graphs that belong to both families.

    \begin{solution}
        Temp
    \end{solution}

    % Problem 2
    \textbf{2.} Prove that removing two opposite corner squares from an
    $8$-by-$8$ checkerboard leaves a subboard that cannot be partitioned
    into $1$-by-$2$ and $2$-by-$1$ rectangles. Using the same argument,
    make a general statement about all bipartite graphs.

    \medskip

    \noindent
    \textit{Hint:} An $8$-by-$8$ checkerboard can be represented by a graph.
    Each square corresponds to a vertex, and two vertices are adjacent if
    and only if the corresponding squares share a side.

    \begin{proof}
        Temp
    \end{proof}

    % Problem 3
    \textbf{3.} Let $G$ be a simple bipartite graph with $n$ vertices and $m$
    edges. Prove that
    \[
        m \leq \frac{n^2}{4}.
    \]

    \medskip

    \noindent
    \textit{Hint:} Suppose $X$ and $Y$ are partite sets of $G$ with
    $|X|=r$ and $|Y|=s$. First show that
    \[
        m \leq rs,
    \]
    and use it to prove
    \[
        m \leq \frac{n^2}{4}.
    \]

    \begin{proof}
        temp
    \end{proof}

    % Problem 4
    \textbf{4.} Let $G$ be a graph with girth $4$ in which every vertex has
    degree $k$. Prove that $G$ has at least $2k$ vertices. Determine all
    such graphs with exactly $2k$ vertices.

    \begin{proof}
        temp
    \end{proof}

    \textbf{5.} Let $P$ be the Petersen graph. Use the definition of $P$ given
    in terms of $2$-element subsets of a set of size $5$ to prove the
    following.

    \begin{enumerate}[label=(\alph*)]
    \item $P$ contains a cycle of length $6$.

    \item $P$ contains a cycle of length $8$.

    \item $P$ does not contain a cycle of length $7$.

    \medskip

    \noindent
    \textit{Hint:} For contradiction, assume that $C$ is a cycle in
    $P$ of length $7$. First, explain why every edge of $P$ whose
    endpoints both lie in $C$ must be an edge of $C$.

    Let $u$ be a vertex of $C$ and let $v,w$ be vertices of $C$
    farthest from $u$. Show that there is some vertex
    \[
        x \notin V(C)
    \]
    such that $u$ and $v$ are both adjacent to $x$, and $u$ and $w$
    are both adjacent to $x$. Conclude that $\{w,v,x\}$ are the
    vertices of a $3$-cycle in $P$, which contradicts the fact that
    $P$ is triangle-free.
    \end{enumerate}

    \begin{solution}
        temp
    \end{solution}

\end{document}